\DocumentMetadata{lang=en-US, pdfversion=2.0, pdfstandard=ua-2, testphase={phase-III,math,title,table,firstaid}}
% CPSY 1291 -- Mean, covariance, eigenvectors: the mathematics of PCA
% Worked examples for the PCA lecture, with the numbers written out and the NumPy beside them.
% Compile:  latexmk -lualatex pca-statistics-notes.tex
\documentclass[11pt]{article}
\usepackage{cpsy1291-handout}
\hypersetup{pdftitle={CPSY 1291: Mean, covariance, eigenvectors: the mathematics of PCA}, pdfauthor={CPSY 1291, Brown University}}

\newcommand{\vx}{\bm{x}}
\newcommand{\vw}{\bm{w}}
\newcommand{\vv}{\bm{v}}
\newcommand{\vz}{\bm{z}}
\newcommand{\vmu}{\bm{\mu}}
\newcommand{\mX}{\bm{X}}
\newcommand{\mC}{\bm{C}}
\newcommand{\mV}{\bm{V}}
\newcommand{\mI}{\bm{I}}
\newcommand{\mL}{\bm{\Lambda}}
\newcommand{\norm}[1]{\left\lVert #1 \right\rVert}
\newcommand{\np}[1]{{\small\color{lightgray}In NumPy: \texttt{#1}}\par}

\begin{document}

\begin{center}
{\sffamily\small\color{accent}CPSY 1291 \ $\cdot$\ Worked examples}\\[5pt]
{\sffamily\LARGE\bfseries\color{brown} Mean, covariance, eigenvectors: the mathematics of PCA}\\[3pt]
{\sffamily\large\color{ink} The statistics and linear algebra behind the PCA lecture and Assignments 1--2, with the numbers written out}\\[6pt]
{\color{lightgray}\small Fall 2026}
\end{center}
\vspace{4pt}\hrule\vspace{10pt}

\section*{Who this is for}
\addcontentsline{toc}{section}{Who this is for}

The PCA lecture paired every formula with a picture and left the algebra to one slide.
This handout does the algebra, one object at a time, on a data set of four points in two
dimensions so that every number can be checked by hand. It goes a little beyond the lecture
in three places: it derives why the top eigenvector is the direction of greatest variance,
it explains where the singular value decomposition comes in, and it works out what
standardization and whitening do to the same four points. The NumPy line after each
definition is the one that Assignment~2's \texttt{fit\_pca} helper uses. Cover each
calculation and try it before reading on; the answers to ``Try it'' are at the end.

\begin{keybox}
\ask{PCA in one sentence.} Centre the data, form the covariance matrix, take its
eigenvectors, sort them by eigenvalue. The eigenvectors are the directions along which the
cloud is spread; the eigenvalues are the variances along those directions.
\end{keybox}

\textbf{The running example.} Four responses to four stimuli, each a vector of two
entries (two neurons, or two pixels):
\[
\vx^{(1)} = (5, 3),\qquad \vx^{(2)} = (1, -1),\qquad \vx^{(3)} = (4, 0),\qquad \vx^{(4)} = (2, 2).
\]
Stacked as rows they form the $4\times2$ response matrix $\mX$, stimuli down, features
across, the convention used throughout the course.

\section{The statistics of one variable}

\subsection{Mean and deviations}

The \emph{mean} of $n$ numbers is their sum divided by $n$. The first feature (the first
column of $\mX$) has values $5, 1, 4, 2$:
\[
\bar{x}_1 = \tfrac{1}{4}(5 + 1 + 4 + 2) = 3, \qquad
\bar{x}_2 = \tfrac{1}{4}(3 - 1 + 0 + 2) = 1 .
\]
The \emph{mean vector} $\vmu = (\bar{x}_1, \bar{x}_2) = (3, 1)$ is the average response,
one point standing in for all four. Subtracting it from every row \emph{centres} the data:
\[
\tilde{\vx}^{(1)} = (2, 2),\qquad \tilde{\vx}^{(2)} = (-2, -2),\qquad
\tilde{\vx}^{(3)} = (1, -1),\qquad \tilde{\vx}^{(4)} = (-1, 1).
\]
The centred rows sum to zero in each column, always: that is what subtracting the mean
does. Everything that follows is computed on $\tilde{\mX}$, and the mean is kept aside so
that it can be added back at the end.

\np{mu = X.mean(axis=0); Xc = X - mu}

\subsection{Variance and standard deviation}

The \emph{variance} of a variable is the mean of its squared deviations from its mean:
\[
\operatorname{var}(x_1) = \tfrac{1}{n}\sum_{i=1}^{n}\big(x^{(i)}_1 - \bar{x}_1\big)^2
= \tfrac{1}{4}(4 + 4 + 1 + 1) = 2.5, \qquad
\operatorname{var}(x_2) = \tfrac{1}{4}(4 + 4 + 1 + 1) = 2.5 .
\]
The \emph{standard deviation} $\sigma$ is its square root, $\sigma_1 = \sigma_2 =
\sqrt{2.5} \approx 1.58$, and has the units of the data; the variance has those units
squared.

\textbf{$n$ or $n-1$.} Software usually divides by $n - 1$ rather than $n$ (NumPy's
\texttt{np.var} divides by $n$ unless told \texttt{ddof=1}; \texttt{np.cov} and
Assignment~2's \texttt{fit\_pca} divide by $n - 1$). The $n - 1$ version is the unbiased
estimate of the variance of the population the sample came from; the $n$ version is the
variance of the sample itself. With $n - 1$ every variance and covariance in this handout
is multiplied by $4/3$, and \emph{nothing else changes}: the eigenvectors are identical and
the eigenvalues all scale by the same factor. The lecture slides use $1/n$ because the
formulas read more simply.

\subsection{z-scores}

A \emph{z-score} rescales a variable so that its mean is $0$ and its standard deviation
is $1$: $z^{(i)} = (x^{(i)} - \bar{x})/\sigma$. For the first feature,
$z_1 = (2, -2, 1, -1)/1.58 = (1.26, -1.26, 0.63, -0.63)$. A z-score of $1.26$ means ``1.26
standard deviations above the mean''. Section~\ref{sec:standardize} shows what happens to
PCA when every feature is z-scored first.

\np{(X - X.mean(0)) / X.std(0)}

\section{Two variables: covariance and correlation}

\subsection{Covariance}

The \emph{covariance} of two variables is the mean product of their deviations:
\[
\operatorname{cov}(x_1, x_2) = \tfrac{1}{n}\sum_{i}\big(x^{(i)}_1 - \bar{x}_1\big)\big(x^{(i)}_2 - \bar{x}_2\big)
= \tfrac{1}{4}(4 + 4 - 1 - 1) = 1.5 ,
\]
the four products being $2\cdot2$, $(-2)(-2)$, $1\cdot(-1)$ and $(-1)\cdot1$.
A positive covariance says that when one variable is above its mean the other tends to be
too; here the two large points $(2,2)$ and $(-2,-2)$ agree in sign and outweigh the two
small ones that disagree. The covariance of a variable with itself is its variance.

\subsection{Correlation}

Covariance has the units of the two variables multiplied. Dividing by both standard
deviations removes the units and gives the \emph{Pearson correlation},
\[
r = \frac{\operatorname{cov}(x_1, x_2)}{\sigma_1\,\sigma_2} = \frac{1.5}{\sqrt{2.5}\sqrt{2.5}} = 0.6,
\]
a number between $-1$ and $1$. Equivalently, $r$ is the covariance of the two z-scored
variables, and (as the first worked-examples handout showed) it is the cosine of the angle
between the two centred columns. On the lecture slide that plotted every pixel against its
neighbour, $r = 0.81$ for the photograph and $-0.01$ for random pixels: that $r$ is this
formula applied to $4{,}032$ pairs.

\np{np.cov(X.T, ddof=0) gives the 2 x 2 matrix; np.corrcoef(X.T) the correlations.}

\subsection{The covariance matrix}
\label{sec:cov}

Put the variances on the diagonal and the covariances off it:
\[
\mC = \begin{pmatrix} \operatorname{var}(x_1) & \operatorname{cov}(x_1,x_2) \\ \operatorname{cov}(x_2,x_1) & \operatorname{var}(x_2) \end{pmatrix}
= \begin{pmatrix} 2.5 & 1.5 \\ 1.5 & 2.5 \end{pmatrix}.
\]
For $d$ features $\mC$ is $d\times d$ and \emph{symmetric}, $C_{jk} = C_{kj}$. The whole
matrix is one product of the centred data with its own transpose,
\[
\mC = \tfrac{1}{n}\,\tilde{\mX}^{\top}\tilde{\mX} = \tfrac{1}{n}\sum_{i=1}^{n}\tilde{\vx}^{(i)}\tilde{\vx}^{(i)\top},
\]
which is the formula on the lecture's eigendecomposition slide. Check the $(1,2)$ entry:
$\tfrac{1}{4}(2\cdot2 + (-2)(-2) + 1\cdot(-1) + (-1)\cdot1) = 1.5$. For a $64\times64$ face
image $d = 4{,}096$ and $\mC$ has sixteen million entries, which is why
Section~\ref{sec:svd} computes PCA without ever forming it.

\np{C = Xc.T @ Xc / len(Xc)  (or np.cov(Xc.T), which divides by n - 1).}

\subsection{The variance along any direction}
\label{sec:quadratic}

The one property of $\mC$ that makes PCA work: for any unit vector $\vw$ ($\norm{\vw} =
1$), the variance of the data \emph{projected onto} $\vw$ is $\vw^{\top}\mC\vw$. The proof
is two lines. The coordinate of point $i$ along $\vw$ is $z^{(i)} = \vw^{\top}\tilde{\vx}^{(i)}$,
a dot product; these coordinates have mean zero because the $\tilde{\vx}^{(i)}$ do; so
\[
\operatorname{var}(z) = \tfrac{1}{n}\sum_i \big(\vw^{\top}\tilde{\vx}^{(i)}\big)^2
= \vw^{\top}\Big(\tfrac{1}{n}\sum_i \tilde{\vx}^{(i)}\tilde{\vx}^{(i)\top}\Big)\vw
= \vw^{\top}\mC\vw .
\]
\textbf{Example.} Along the first axis, $\vw = (1, 0)$: $\vw^{\top}\mC\vw = 2.5$, the
variance of $x_1$, as it should be. Along the diagonal, $\vw = (1, 1)/\sqrt{2}$:
\[
\vw^{\top}\mC\vw = \tfrac{1}{2}\,(1, 1)\begin{pmatrix} 2.5 & 1.5 \\ 1.5 & 2.5 \end{pmatrix}\begin{pmatrix}1\\1\end{pmatrix}
= \tfrac{1}{2}(1, 1)\cdot(4, 4) = 4 .
\]
Along the anti-diagonal, $\vw = (1, -1)/\sqrt{2}$: $\tfrac{1}{2}(1,-1)\cdot(1,-1) = 1$.
Check the first one directly: the coordinates along $(1,1)/\sqrt 2$ are
$z = (4, -4, 0, 0)/\sqrt{2} = (2.83, -2.83, 0, 0)$, whose mean square is $(8 + 8)/4 = 4$.
The cloud is four times as spread along the diagonal as across it.

\textit{Try it.} Compute $\vw^{\top}\mC\vw$ for $\vw = (0, 1)$ and for $\vw = (0.6, 0.8)$.

\np{w @ C @ w, with w a unit vector; np.var(Xc @ w).}

\section{Eigenvectors and eigenvalues}

\subsection{Definition, and the 2 x 2 case by hand}

A vector $\vv \neq \bm{0}$ is an \emph{eigenvector} of a square matrix $\mC$ if
multiplying by $\mC$ only rescales it:
\[
\mC\vv = \lambda\vv ,
\]
and the scale factor $\lambda$ is its \emph{eigenvalue}. Most vectors are turned by $\mC$;
eigenvectors are the ones that keep their direction.

For a $2\times2$ matrix the eigenvalues can be found by hand. $\mC\vv = \lambda\vv$ means
$(\mC - \lambda\mI)\vv = \bm{0}$ with $\vv \neq \bm 0$, which is only possible when the
matrix $\mC - \lambda\mI$ squashes some direction to zero, that is, when its determinant is
zero:
\[
\det\begin{pmatrix} 2.5-\lambda & 1.5 \\ 1.5 & 2.5-\lambda \end{pmatrix}
= (2.5-\lambda)^2 - 1.5^2 = 0
\quad\Longrightarrow\quad 2.5 - \lambda = \pm 1.5
\quad\Longrightarrow\quad \lambda_1 = 4,\ \lambda_2 = 1 .
\]
For $\lambda_1 = 4$, solve $(\mC - 4\mI)\vv = \bm 0$: the first row reads
$-1.5\,v_1 + 1.5\,v_2 = 0$, so $v_1 = v_2$ and the unit eigenvector is
$\vv_1 = (1, 1)/\sqrt{2}$. For $\lambda_2 = 1$ the first row reads
$1.5\,v_1 + 1.5\,v_2 = 0$, so $\vv_2 = (1, -1)/\sqrt{2}$. Check:
\[
\mC\vv_1 = \tfrac{1}{\sqrt2}\begin{pmatrix} 2.5 + 1.5 \\ 1.5 + 2.5\end{pmatrix}
= \tfrac{1}{\sqrt2}\begin{pmatrix}4\\4\end{pmatrix} = 4\,\vv_1, \qquad
\mC\vv_2 = \tfrac{1}{\sqrt2}\begin{pmatrix} 2.5 - 1.5 \\ 1.5 - 2.5\end{pmatrix} = 1\cdot\vv_2 .
\]
These are exactly the two directions of Section~\ref{sec:quadratic}, and the eigenvalues
are exactly the variances found there: $4$ along the diagonal, $1$ across it. That is not a
coincidence, and the next subsection says why.

Beyond $2\times2$ nobody solves the polynomial; the computer does it numerically. Two things
to know about the call: \texttt{np.linalg.eigh} is for symmetric matrices (the \texttt{h}
is for Hermitian) and is faster and more accurate than the general \texttt{eig}, and it
returns the eigenvalues in \emph{ascending} order, so the columns have to be reversed to
put the largest first. An eigenvector is only defined up to sign: $-\vv_1$ satisfies the
same equation, and different libraries, or the same library on different machines, may
return either. Assignment~2's helper fixes the sign so that the largest entry of each
eigenvector is positive.

\np{lam, V = np.linalg.eigh(C); lam, V = lam[::-1], V[:, ::-1]}

\subsection{Why the eigenvectors of C are the principal axes}
\label{sec:why}

Two facts about symmetric matrices, stated without proof: their eigenvalues are real, and
their eigenvectors can be chosen \emph{orthonormal}, unit length and mutually
perpendicular. (Check: $\vv_1\cdot\vv_2 = \tfrac12(1\cdot1 + 1\cdot(-1)) = 0$.) So the
eigenvectors form a basis, and any unit vector can be written in it,
$\vw = a_1\vv_1 + a_2\vv_2 + \dots + a_d\vv_d$ with $\sum_k a_k^2 = 1$. Now compute the
variance along $\vw$ using $\mC\vv_k = \lambda_k\vv_k$ and the orthonormality:
\[
\vw^{\top}\mC\vw = \Big(\sum_j a_j\vv_j\Big)^{\top}\Big(\sum_k a_k\lambda_k\vv_k\Big)
= \sum_k \lambda_k a_k^2 .
\]
This is a weighted average of the eigenvalues with non-negative weights $a_k^2$ that sum
to one. Its largest possible value is the largest eigenvalue, $\lambda_1$, reached when all
the weight sits on $\vv_1$ ($a_1 = \pm1$, the rest zero); its smallest is $\lambda_d$.
\emph{The direction of greatest variance is the top eigenvector, and the variance along it
is the top eigenvalue.} Among directions perpendicular to $\vv_1$ the same argument gives
$\vv_2$, and so on. That is the whole justification of PCA, in four lines.

\textbf{Example.} $\vw = (0.6, 0.8)$ has $a_1 = \vw\cdot\vv_1 = 1.4/\sqrt2 = 0.99$ and
$a_2 = \vw\cdot\vv_2 = -0.2/\sqrt2 = -0.14$, so
$\vw^{\top}\mC\vw = 4\,(0.99)^2 + 1\,(0.14)^2 = 3.92 + 0.02 = 3.94$, which you can confirm
by the direct product.

\subsection{Total variance, and the Pythagorean split}

The sum of the variances of all the features, $\operatorname{tr}\mC = \sum_j C_{jj}$, is
the \emph{total variance}; it is also the mean squared distance of the centred points from
the origin, $\tfrac1n\sum_i\norm{\tilde{\vx}^{(i)}}^2$. It equals the sum of the
eigenvalues (a fact about every square matrix): here $2.5 + 2.5 = 4 + 1 = 5$.

Project every point onto $\vw$ and it splits into a part along $\vw$ and a perpendicular
remainder. By Pythagoras the squared lengths add, and averaging over the cloud,
\[
\underbrace{\tfrac1n\sum_i\norm{\tilde{\vx}^{(i)}}^2}_{\text{total, } 5}
= \underbrace{\vw^{\top}\mC\vw}_{\text{projected variance}}
+ \underbrace{\tfrac1n\sum_i\norm{\tilde{\vx}^{(i)} - z^{(i)}\vw}^2}_{\text{reconstruction error}} .
\]
The left side does not depend on $\vw$. So whichever $\vw$ captures the most variance
leaves the least error: \emph{maximising projected variance and minimising reconstruction
error are the same problem}, and $\vv_1$ solves both. For $\vw = \vv_1$ the projected
variance is $4$ and the error is $1 = \lambda_2$; for $\vw = (1, 0)$ they are $2.5$ and
$2.5$.

\section{Projection and reconstruction}
\label{sec:proj}

\subsection{Coordinates in the new basis}

Collect the top $K$ eigenvectors as the columns of $\mV_K$ ($d\times K$). The
\emph{coordinates} of a point in the new basis, its PCA scores, are the $K$ dot products
\[
\vz = \mV_K^{\top}(\vx - \vmu), \qquad z_k = \vv_k^{\top}(\vx - \vmu).
\]
With $K = d$ this is a pure rotation: nothing is lost. With $K = 1$ and our four points,
$z = \vv_1^{\top}\tilde{\vx}$ gives $(2.83, -2.83, 0, 0)$: two points far out along the
diagonal, two at the origin. Points 3 and 4, which differ from each other by $(2, -2)$,
receive the same single coordinate, because that difference lies entirely along $\vv_2$.

\subsection{The way back, and what is lost}

Reconstruction reverses the projection and adds the mean back:
\[
\hat{\vx} = \vmu + \mV_K\vz = \vmu + \sum_{k=1}^{K} z_k\vv_k .
\]
With $K = 1$: $\hat{\vx}^{(1)} = (3, 1) + 2.83\,(1, 1)/\sqrt2 = (3, 1) + (2, 2) = (5, 3)$,
exact, because point 1 lay on the $\vv_1$ line to begin with. But
$\hat{\vx}^{(3)} = (3, 1) + 0 = (3, 1)$, while the true point is $(4, 0)$: the error
vector is $(1, -1)$, of squared length $2$. Point 4's error is $(-1, 1)$, also $2$; points
1 and 2 have none. Mean squared error over the four points: $(0 + 0 + 2 + 2)/4 = 1 =
\lambda_2$. In general the mean squared reconstruction error with $K$ components is the
sum of the discarded eigenvalues, $\sum_{k > K}\lambda_k$, which is what the lecture's
``residual fraction'' curves plot: the fraction $\sum_{k>K}\lambda_k / \sum_k\lambda_k$,
here $1/5 = 20\%$ for $K=1$. The complementary number, $\lambda_1 / \sum_k\lambda_k =
80\%$, is the ``variance explained'' of the scree plot.

\np{Z = Xc @ V[:, :K]; X\_hat = mu + Z @ V[:, :K].T}

\textit{Try it.} Reconstruct point 2 with $K = 1$ and with $K = 2$.

\section{Computing PCA without forming C: the SVD}
\label{sec:svd}

Every matrix has a \emph{singular value decomposition}, $\tilde{\mX} = \bm{U}\bm{S}\mV^{\top}$,
with $\bm{U}$ ($n\times r$) and $\mV$ ($d\times r$) having orthonormal columns and $\bm S$
a diagonal matrix of non-negative \emph{singular values} $s_1 \ge s_2 \ge \dots$, where $r$
is the rank of $\tilde{\mX}$. Substitute it into the covariance:
\[
\mC = \tfrac1n\tilde{\mX}^{\top}\tilde{\mX} = \tfrac1n\,\mV\bm S\bm U^{\top}\bm U\bm S\mV^{\top}
= \mV\Big(\tfrac{1}{n}\bm S^2\Big)\mV^{\top},
\]
using $\bm U^{\top}\bm U = \mI$. Compare with the eigendecomposition $\mC = \mV\mL\mV^{\top}$:
the right singular vectors of the centred data \emph{are} the eigenvectors of $\mC$, and
$\lambda_k = s_k^2/n$ (or $s_k^2/(n-1)$ with the other convention). This is how
scikit-learn and Assignment~2's \texttt{fit\_pca} compute PCA: the SVD of the $n\times d$
data never builds the $d\times d$ covariance, which for images would not fit in memory.

The rank is at most $\min(n - 1, d)$: the $-1$ is because centring removes one degree of
freedom (the rows sum to zero). With four points in two dimensions, $r \le 2$; with $400$
faces of $4{,}096$ pixels, at most $399$ eigenvalues are non-zero however many pixels
there are, a fact that surprises people the first time a scree plot stops early.

\np{U, s, Vt = np.linalg.svd(Xc, full\_matrices=False); lam = s**2 / len(Xc); V = Vt.T}

\section{Standardization: the same points, different units}
\label{sec:standardize}

PCA maximises variance, and variance depends on units. Suppose the first feature had been
recorded in different units, ten times larger: $x_1 = 50, 10, 40, 20$. Then
$\operatorname{var}(x_1) = 250$, $\operatorname{cov}(x_1, x_2) = 15$, and
\[
\mC' = \begin{pmatrix} 250 & 15 \\ 15 & 2.5 \end{pmatrix}, \qquad
\lambda_1 \approx 250.9,\ \vv_1 \approx (0.998, 0.060) .
\]
The first component now points almost exactly along the first feature and carries $99\%$
of the variance: the second feature has become invisible, although the \emph{shape} of the
cloud, in the sense of which points are near which, has not changed at all. (Check the
eigenvalues quickly: $\lambda_1 + \lambda_2 = 252.5$ and $\lambda_1\lambda_2 = \det\mC' =
625 - 225 = 400$.)

Z-scoring each feature first (Section 1.3) gives every feature variance $1$, so the
covariance matrix of the z-scores is the \emph{correlation matrix},
\[
\mC_z = \begin{pmatrix} 1 & r \\ r & 1 \end{pmatrix} = \begin{pmatrix} 1 & 0.6 \\ 0.6 & 1\end{pmatrix},
\qquad \lambda_1 = 1.6,\ \lambda_2 = 0.4,\ \vv_1 = (1, 1)/\sqrt 2 ,
\]
whatever the original units were (for a $2\times2$ correlation matrix the eigenvalues are
always $1 \pm r$). Standardize when features are in different units or have variances that
differ for uninteresting reasons; do not standardize when the differences in variance are
the signal, as when every feature is a neuron's firing rate and some neurons really are
more active. Assignment~1 standardizes the IT units before PCA for the first reason.

\np{Xz = (X - X.mean(0)) / X.std(0); then PCA as before.}

\section{Whitening: one step further}

Projecting onto all $d$ eigenvectors rotates the cloud so that its coordinates are
\emph{uncorrelated}: the covariance of $\vz = \mV^{\top}\tilde{\vx}$ is $\mV^{\top}\mC\mV =
\mL$, diagonal, with the eigenvalues down the diagonal. Dividing coordinate $k$ by
$\sqrt{\lambda_k}$ then gives every coordinate variance $1$:
\[
\vz_w = \mL^{-1/2}\mV^{\top}(\vx - \vmu), \qquad \operatorname{cov}(\vz_w) = \mI .
\]
This is \emph{whitening}: the cloud becomes a round ball of unit radius, every direction
equally spread. On the four points, the whitened coordinates are
$z_1/\sqrt4 = (1.41, -1.41, 0, 0)$ and $z_2/\sqrt1 = (0, 0, 1.41, -1.41)$, and both have
mean square $1$. Whitening is what a learner wants as input when no direction should be
privileged; it is also, to a first approximation, what the retina does to the
correlations between neighbouring pixels. It is not the same as z-scoring: z-scoring
rescales the original features, whitening rescales the principal coordinates, and only the
second removes correlations.

\np{Zw = (Xc @ V) / np.sqrt(lam)}

\section{All of it in code}

The following function is, line for line, the PCA of Assignment~2 (its \texttt{fit\_pca}
helper divides by $n-1$ and fixes eigenvector signs; both are included). Run it on the
four points and every number in this handout comes back, multiplied by $4/3$ where the
convention says so.

\begin{tcolorbox}[colback=codebg, colframe=codebg, arc=0pt, boxrule=0pt, left=4pt, right=4pt, top=2pt, bottom=2pt]
\small\begin{verbatim}
import numpy as np

X = np.array([[5, 3], [1, -1], [4, 0], [2, 2]], dtype=float)  # (n, d)

def fit_pca(X):
    mean = X.mean(axis=0)                             # section 1.1
    centred = X - mean
    _, s, Vt = np.linalg.svd(centred, full_matrices=False)  # section 6
    eigenvalues = s**2 / (len(X) - 1)                 # n - 1 convention
    eigenvectors = Vt.T                               # one unit eigenvector per column
    for j in range(eigenvectors.shape[1]):            # deterministic sign
        if eigenvectors[np.argmax(np.abs(eigenvectors[:, j])), j] < 0:
            eigenvectors[:, j] *= -1
    return mean, eigenvalues, eigenvectors

mean, lam, V = fit_pca(X)
print(lam)                      # [5.333 1.333]  = (4, 1) * 4/3
print(V[:, 0])                  # [0.707 0.707]  = v1
Z = (X - mean) @ V[:, :1]       # section 5.1: one coordinate per point
X_hat = mean + Z @ V[:, :1].T   # section 5.2: reconstruction with K = 1
print(((X - X_hat)**2).sum(1))  # [0. 0. 2. 2.]: residuals, point by point
print(lam[1:].sum() / lam.sum())  # 0.2: discarded fraction, either convention
\end{verbatim}
\end{tcolorbox}

Two habits worth keeping from this. Never compare eigenvalues across two analyses without
checking they used the same $n$ or $n-1$ convention, and never compare eigenvector signs
across two analyses at all; compare the variance fractions and the absolute values of the
projections instead.

\subsection*{Answers to ``Try it''}
\S\ref{sec:quadratic}: $\vw = (0, 1)$ gives $2.5$, the variance of $x_2$;
$\vw = (0.6, 0.8)$ gives $0.36\cdot2.5 + 2\cdot0.48\cdot1.5 + 0.64\cdot2.5 = 0.9 + 1.44 + 1.6 = 3.94$,
matching \S\ref{sec:why}.\quad
\S\ref{sec:proj}: point 2 has $z_1 = -2.83$ and $z_2 = 0$, so
$\hat{\vx}^{(2)} = (3, 1) - 2.83\,(1,1)/\sqrt2 = (1, -1)$ already with $K = 1$; $K = 2$ adds
$0\cdot\vv_2$ and changes nothing. Points 3 and 4 are the ones that need the second
component.

\subsection*{Symbols used in this handout}
\begin{center}
\footnotesize
\setlength{\tabcolsep}{4pt}
\begin{tabular}{@{}ll@{\quad}ll@{}}
\toprule
$\mX$, $\tilde{\mX}$ & $n\times d$ data; the same, centred & $\mC$ & $d\times d$ covariance matrix, $\tfrac1n\tilde{\mX}^{\top}\tilde{\mX}$ \\
$\vmu$, $\bar{x}_j$ & mean vector; mean of feature $j$ & $\lambda_k$, $\vv_k$ & $k$-th eigenvalue and unit eigenvector of $\mC$ \\
$\sigma_j$, var, cov & standard deviation; variance; covariance & $\mV_K$, $\mL$ & top $K$ eigenvectors; diagonal of eigenvalues \\
$r$ & Pearson correlation, $\operatorname{cov}/(\sigma_1\sigma_2)$ & $\vz$, $z_k$ & PCA coordinates (scores); the $k$-th one \\
$\vw$, $\vw^{\top}\mC\vw$ & a unit direction; the variance along it & $\hat{\vx}$ & reconstruction, $\vmu + \mV_K\vz$ \\
$\operatorname{tr}\mC$ & total variance, $\sum_j C_{jj} = \sum_k\lambda_k$ & $\bm U\bm S\mV^{\top}$, $s_k$ & the SVD; singular values \\
\bottomrule
\end{tabular}
\end{center}

\end{document}
